% Language simplified for clearer student-style academic writing; cited claims are unchanged.
\section{Preliminaries}

A unified cattle-monitoring system needs to connect different kinds of visual information. Body condition is mainly related to body shape and morphology, behavior is related to posture and movement, and identity depends on features that make one animal different from another. Because of this, the thesis connects ideas from livestock monitoring, visual representation learning, and multi-task learning. This chapter reviews how earlier methods select and use visual information, and how they are evaluated under different conditions.

The chapter first explains the main concepts needed for the three tasks. It then reviews cattle-specific models, foreground processing, anatomy and pose, temporal information, and task-specific feature sharing. These topics are brought together to explain the research gap of this thesis.

\subsection{Visual Representations and Transfer Learning}

A visual representation is the set of features a model uses to describe an image or video sequence. Older recognition systems often depended on hand-designed measurements. Deep neural networks can instead learn useful features directly from examples. In a convolutional network, simple local patterns are gradually combined into higher-level features. The features passed to the final prediction head decide what information is available for making the prediction.

Reusable image encoders have made this easier in specialized fields. ResNet uses residual connections to make deep networks easier to train~\cite{he2016resnet}. EfficientNet focuses on scaling network depth, width, and input resolution together~\cite{efficientnet2019}. These models solve different design problems, but both are useful starting points for transfer learning. Instead of learning every visual feature from a small cattle dataset, a pretrained model can be adapted to the cattle task.

Transfer learning and cattle-centered representation learning are not the same in terms of functionality. Transfer learning focuses on the source of initial weights of the model. Cattle-centered learning gives us a clue as to what kind of input data was fed to the model. The same pre-trained encoder may be used with different types of input data such as full image, crop of the cow, masked cow, and anatomically detailed images. That is why the encoder will stay the same regardless of the changes in the input data. This difference is important when comparing why one representation performs better than another.

\subsection{Prediction Tasks and Their Visual Requirements}

There are different approaches for predicting BCS, which are classification, ordinal prediction, and regression. In classification, all the scores can be treated as individual classes. In ordinal prediction, order is significant since the prediction of a neighboring score is a less significant error compared to the one which is distant.  However, regression just tries to predict the numeric output. CORAL is one example of an algorithm designed for the ordered labels~\cite{coral2020}. This algorithm may teach the classifier about the order of BCS, but it will not choose the visual features of the cow.

In the case of behavior recognition, a class is actually an activity that is not measured by numbers or scores. One frame can be able to inform us about the posture of a cow , whereas multiple frames can show how the cow moves over time. It is important since there are activities that might be quite similar in one frame only, for example, standing and walking. which can look almost the same in a single frame. So, the model should use a suitable amount of time for the detection task.

Identification and Re-ID are two distinct tasks. Closed set classifier identifies the cows seen in the training dataset. On the other hand , the query image is compared with the gallery images based on the embedding learned in Re-ID. It becomes possible for the system to be tested on cows unseen during the training phase. The good Re-ID embedding will classify images of the same cow into one cluster and images of the other cows into different clusters. The significance of this technique in cattle research lies in the changing herd of cattle~\cite{andrew2021opencows}.

\subsection{Multi-Task Learning and Negative Transfer}

Multi-task learning entails the training of different tasks in the same model where particular parts of the model are common and the training objective is minimizing the total loss of all the tasks,
\begin{equation}
    \mathcal{L}_{\mathrm{MTL}}=\sum_{t=1}^{T}\lambda_t\mathcal{L}_t,
\end{equation}
where in this equation, $\mathcal{L}_t$ is the loss of task $t$, and $\lambda_t$ determines the weight of the loss of task $t$. This equation calculates the total loss; however, but there is no information about what amount of the network should be shared. A neural network model can share all layers, a few layers, or exchange information between tasks.

In the survey paper by Crawshaw~\cite{crawshaw2020mtl}, MTL is considered in the context of information sharing, task-specific components, and optimization. MTL could be advantageous if one task is providing some useful information for another task. However, joint training could be harmful to a specific task; that is, the effect is referred to as negative transfer. It might occur due to the inappropriateness of information shared between tasks, conflicting task gradients, and unbalanced training losses. It is critical to distinguish information sharing and training for the purpose of the thesis.

\section{Review of Existing Research}

\subsection{From Automated Observations to Integrated Livestock Monitoring}

Livestock-monitoring research has moved from separate measurements toward systems that combine sensing, prediction, and decision support. This can be seen in both computer-vision reviews and broader precision-livestock studies. Antognoli et al.~\cite{antognoli2025ani1} review computer vision in dairy management, while Zin and Tin~\cite{zin2026ab26} discuss computer vision across wider livestock applications. These reviews show that image recognition is only one part of monitoring. The image first needs to be turned into a reliable observation about the animal before it can support a farm decision.

Research on livestock data analysis also shows why integration is important. Palma et al.~\cite{palma2025ani1} report that predictive analysis is much more common than prescriptive decision support in the dairy studies they review. Sani et al.~\cite{sani2026ab25} also discuss the difficulty of moving precision livestock technologies into practical use. This means that collecting more data or improving one classifier is not enough. The outputs also need to be connected over time and related to the correct animal and production setting.

Sensor-based and vision-based systems provide different information. Lee et al.~\cite{lee2026ab26} discuss how biosensors can provide continuous measurements at the individual-animal level, while cameras directly record visible appearance and activity. Vision is especially useful for this thesis because body shape, posture, movement, and coat patterns are all visible in images or video. This gives the three tasks a common sensing source, even though they need different information from it.

This does not mean that all three tasks should use exactly the same features. A unified system can share cameras, preprocessing, and some learned features while keeping task-specific processing where needed. Therefore, integration is not only about putting several outputs in one system. It is also about deciding what information can be shared without losing the needs of each task.

\subsection{Body Condition Scoring: Anatomy, Image Geometry, and Ordered Prediction}

Automated BCS has a strong anatomical basis. The scoring system introduced by Edmonson et al.~\cite{edmonson1989bcs} and the descriptors studied by Ferguson et al.~\cite{ferguson1994bcs} connect BCS values to the appearance of specific body regions. For computer vision, this means that overall body size is not enough. The model needs information about tissue coverage, body contours, and visible anatomical structures. The camera view and the selected body region therefore matter.

One approach uses depth images to capture body geometry more directly. Rodr\'iguez Alvarez et al.~\cite{rodriguez2018bcs} use convolutional networks on depth-based cattle images and report results using clear BCS error tolerances. Their work combines geometric sensing with learned prediction instead of depending only on hand-designed body measurements. It also shows why exact accuracy and tolerance-based accuracy should be reported separately. BCS has an ordered scale, so a small error near the correct score is different from a much larger error.

Another approach looks for useful condition cues in normal RGB images. Liu et al.~\cite{liu2025vets} focus on tail-related features and use EfficientNet-B0 with channel and spatial attention. Their method combines a selected anatomical region with an efficient model, and uses distillation to reduce the model size. This shows that a BCS model can be specialized both by where it looks and by the network design. It also differs from simply using a larger backbone on a less controlled image region.

Viewpoint also changes what information is available. Yao et al.~\cite{yao2026jds2} use localized side-view cattle images with a ConvNeXt-based regression model and cow-level cross-validation. A tail-focused method sees a small anatomical region, while a side-view method sees a larger body profile with different geometry. These are not only different model choices; they expose different parts of the animal and can therefore provide different BCS information.

Research on cattle body measurement also shows why image size and real body geometry should be separated. PickAMoo, developed by Guzhva et al.~\cite{guzhva2026s415}, combines smartphone images, LiDAR-based scaling, and Mask R-CNN segmentation for cattle weight estimation. Weight is not the same target as BCS, but the method is relevant because it measures the animal after isolating it and links image size to physical scale. This supports a simple design idea: a cow should not be judged as having a different biological condition just because it appears larger in the image.

Together, these studies show three important design choices for automated BCS. First, the image must contain useful condition information, including the right anatomy and viewpoint. Second, the representation must extract that information through a crop, contour, depth image, attention mechanism, or another suitable method. Third, the prediction method must match the ordered target. CORAL handles the order between output classes~\cite{coral2020}, but it still depends on the input containing useful condition information. In this thesis, preserving morphology is therefore a representation goal, while ordinal prediction is a separate modeling choice.

\subsection{Behavior Recognition: Posture, Motion, and Environmental Context}

Behavior classes can differ in both appearance and duration. Lying may often be recognized from body posture, while walking is better described by movement. Feeding and drinking also involve interaction with a resource such as a trough. Because of this, behavior research uses different amounts of spatial and temporal information, and a single-frame classifier cannot describe every type of activity equally well.

Image-based behavior recognition shows that posture can already provide useful information. Asim et al.~\cite{asim2026s260} use overhead-camera images and object-detection models to recognize activities including standing, lying, grazing, and estrus-related behavior. Their method is vision-based rather than wearable-sensor based. It shows one way to connect animal localization and activity labels in a frame. A sequence model can then add information about how those visible postures change over time.

Video datasets extend the task beyond one frame. CBVD-5, introduced by Li et al.~\cite{li2024cbvd5}, contains standing, lying, foraging, rumination, and drinking. CVB, introduced by Zia et al.~\cite{zia2023cvb}, provides densely sampled behavior clips from an open-pasture setting and includes a SlowFast benchmark. The datasets have different environments and behavior labels, but both show that behavior recognition depends on the animal and on how the activity is defined and observed over time.

MmCows adds synchronized sensors and individual-animal information to video data~\cite{vu2024mmcows}. Its annotations include identity, location, and behavior across several camera views. This allows researchers to ask whether behavior predictions stay consistent across views and animals. A dense video dataset with only temporary tracking IDs can still provide useful motion examples, but it does not provide the same level of biological identity information.

This shows two separate properties of behavior data: \emph{temporal detail} and \emph{identity continuity}. Dense video supports short-term motion analysis. Reliable cow IDs support testing on different animals and following animals over time. One property cannot replace the other. Therefore, a study that cares about both movement and generalization should consider both instead of choosing a dataset only because it contains many images.

Background context can also help and hurt. A trough or feeding barrier may help the model understand feeding or drinking, but the model may also learn that one activity usually happens in one camera area. Removing the whole background can reduce that shortcut but may also remove useful interaction information. The important question is not whether context is always good or bad. The question is how much context is useful for understanding the activity and whether that information still works when the scene changes.

This connects behavior recognition to the cattle-centered idea of the thesis. The cow should remain the main source of visual information, while posture, motion, and nearby interaction context can still be used when needed. It also supports class-sensitive evaluation. A model may perform well on common resting behaviors and still fail on rarer or more dynamic behaviors. Overall accuracy alone may hide this problem, especially when class frequency is linked to the recording source.

\subsection{Individual Identification and Re-ID: From Coat Patterns to Cross-Setting Recognition}

Visual cattle identification depends on features that are different between animals. Andrew et al.~\cite{andrew2017cattle} use local coat-pattern matching in RGB-D images. Their method makes the identity information clear: local markings can help distinguish Holstein-Friesian cows. It also shows an important limitation that still applies to modern models. The useful identity feature must be visible in both images being compared.

Deep metric learning changes how these comparisons are learned. Andrew et al.~\cite{andrew2021opencows} use an embedding-based method with a combined classification and reciprocal-triplet objective. Instead of only using a classifier for known cows, the embedding supports similarity-based matching. This allows the feature space to be tested on cows that were not used to train the encoder. The goal therefore changes from memorizing fixed identity labels to learning features that can separate animals more generally.

Face-based identification uses a different body region. Weng et al.~\cite{weng2022cattleface} use a two-branch convolutional network to combine cattle-face information. Compared with coat-based methods, this focuses on facial differences between cows. Which approach is more useful depends on the camera view. A clear face may be hard to obtain from an overhead or rear view, while a useful coat pattern may be hidden by occlusion or missing from a tight face crop. Camera position and visible body coverage are therefore important parts of Re-ID.

Another line of work reduces the amount of manual identity labeling needed in video. Gao et al.~\cite{gao2021cows2021} use the Cows2021 dataset to study self-supervision from tracking-based observations and normalized cow views. Yu et al.~\cite{yu2025multicam} extend this idea to several cameras on a working farm with MultiCamCows2024. These methods use the continuity of video to connect images that are likely to show the same cow. A short tracklet can therefore help the model learn appearance consistency without requiring a manually assigned cow ID for every frame.

Tracking and Re-ID are still different tasks. Tracking keeps an animal associated while it remains continuously visible. Re-ID has to recover that association after a break in time or a change in viewing conditions. Tracking information can help train a Re-ID model, but a temporary track ID should not be treated as a permanent biological cow ID. Multi-camera and cross-setting evaluation are useful because they test larger changes than a continuous video sequence.

Long-term recognition adds another challenge. The BECA dataset of Zhang et al.~\cite{zhang2026beca} includes a large-population beef-cattle setting and a longitudinal setting. This goes beyond short-term matching of cows with distinctive dairy coat patterns. An animal's appearance can change with growth, posture, and the visible body region, and a larger herd creates more visually similar animals. A Re-ID model should therefore be tested not only on the cows used during training, but also on whether its identity features remain useful across larger populations and longer time periods.

Foreground and anatomy processing must also protect the information needed for identity. Removing the background can push the model to focus on the cow, but reducing the cow to only a silhouette or skeleton can remove important coat markings. Keeping the entire scene can create the opposite problem, where the model learns the cow's usual location. Re-ID therefore needs selective representation: reduce environmental shortcuts while keeping the cow-specific visual details needed for matching.

\subsection{Localization and Segmentation as Representation Design}

Localization decides which cow the model should analyze. This is especially important when several cows appear in the same image, because the predicted behavior or identity must belong to the correct animal. A bounding box gives a practical region of interest while still keeping some surrounding context. Cropping can reduce scale differences and focus the encoder, but it also changes how much body detail and background remain in the image.

Modern detectors provide different ways to find this region. RT-DETR, introduced by Zhao et al.~\cite{zhao2024rtdetr}, uses an end-to-end transformer detector designed around the speed--accuracy trade-off. Instance-segmentation methods provide more detail than a box. Mask R-CNN adds a mask branch to the detection framework so that it predicts both the object location and an object-specific foreground mask~\cite{he2017maskrcnn}. In simple terms, a box tells us where the cow is, while a mask estimates which pixels belong to it.

Promptable segmentation separates target selection from mask generation. Segment Anything, proposed by Kirillov et al.~\cite{kirillov2023sam}, can create a mask from prompts such as points or boxes. SAM 2 extends this idea to images and video using a memory-based design~\cite{ravi2024sam2}. In a cattle pipeline, this means that a detector can first find the cow and the segmentation model can then create a candidate mask without training a new cattle-specific segmentation network from scratch. The final mask still depends on how the target was selected and prompted.

Cattle-specific datasets provide more direct supervision. CattleEyeView includes top-down annotations for detection, counting, tracking, pose, and instance segmentation~\cite{ong2023cattleeyeview}. Because several types of annotation describe the same scenes, it is useful for studying cattle perception. However, these perception tasks are not the same as predicting BCS, behavior, and identity. For this thesis, they are useful ways of describing the cow, not proof that one description should be shared by every downstream task.

Recent cattle recognition work also links foreground estimation with identity learning. CATR-Net combines segmentation and cattle-face recognition using edge-based segmentation and attention-based recognition components~\cite{feng2025catr}. This directly connects target separation with identity features. Together with earlier coat- and face-based methods, it shows that cattle-centered recognition is already an active research area. The open question for a broader monitoring system is whether foreground emphasis also helps tasks that depend more on morphology or motion.

Segmentation and feature attention are also different. CBAM changes the importance of channel and spatial features inside a network~\cite{cbam2018}, while a foreground mask changes or weights the image content using an estimated object boundary. Attention can focus on useful regions while still keeping the full image. Masking can reduce background information before the main encoder. The two methods may work together, but they are not the same intervention.

Masking can also be used at different strengths. A hard foreground image removes everything outside the estimated boundary. A softer mask reduces background influence while keeping some context and allowing for uncertain boundaries. The best choice depends on the task. BCS needs useful body contours, behavior may need nearby interaction context, and Re-ID needs coat details. Therefore, mask quality should be judged together with downstream task performance, not only by how clean the mask looks.

\subsection{Pose, Anatomy, and Viewpoint}

Pose estimation represents an animal using the positions of selected body parts. DeepLabCut shows that user-defined landmarks can be learned with transfer learning from a relatively small amount of annotation~\cite{mathis2018deeplabcut}. This changes the input from dense appearance information to a structured description of body configuration. Such information can support movement analysis and can also make some parts of the model input easier to understand.

SuperAnimal extends this idea by combining different animal-pose datasets into transferable models~\cite{ye2024superanimal}. Its training method handles differences between landmark definitions and supports both zero-shot use and further adaptation. This is useful for cattle work because it reduces the need to create a completely new pose dataset from the beginning. CattleEyeView provides a more cattle-specific option in a top-down setting~\cite{ong2023cattleeyeview}. Together, these works show two ways to obtain structural information: transfer from many animal datasets or training with cattle-specific annotations.

The landmark definition matters as much as the pose model itself. A skeleton based on the head, torso, and leg joints can describe posture, but BCS also depends on visible anatomy and tissue coverage. The BCS descriptors studied by Ferguson et al.~\cite{ferguson1994bcs} therefore require more than joint locations. Pose and anatomy are related, but they are not the same representation. BCS may need structural coordinates together with visual detail, while behavior may depend more directly on how the coordinates change.

Viewpoint decides which body parts are visible and how the cow appears in the image. Rear, side, and overhead views can show different body contours even when the cow's real condition has not changed. This matters for Re-ID because the cow should still be matched across views, and for BCS because important regions may be hidden. Viewpoint means the orientation of the animal relative to the camera; it is not the same as camera identity.

The MOO preprint by Grolleau et al.~\cite{grolleau2026moo} studies viewpoint using synthetic cattle rendered from controlled directions. Its angle labels make it possible to study geometric changes separately from some uncontrolled changes in natural images. The authors also study transfer to real cattle images. Synthetic data therefore provide controlled viewpoint labels, while real images test whether the learned information still works under natural appearance and environmental changes.

The main idea is that invariance should be selective. A Re-ID representation may need to stay stable when the same cow is seen from another angle, but a behavior representation should still react when the posture changes. BCS may need stable body-shape information even when image scale changes. Pose and viewpoint can help describe these changes, but the useful response depends on the task. This is why the thesis studies task-conditioned structural information instead of forcing one invariant representation on every output.

\begin{table}[tbp]
\centering
\caption[Visual information and task requirements]{Visual information and task requirements synthesized from the reviewed literature. The entries describe design considerations, not measured outcomes of the proposed system.}
\label{tab:ch2_information_requirements}
\small
\begin{tabularx}{\textwidth}{@{}>{\raggedright\arraybackslash}p{2.2cm}YYY@{}}
\toprule
\textbf{Visual cue} & \textbf{BCS} & \textbf{Behavior} & \textbf{Re-ID}\\
\midrule
Foreground appearance & Body contours and local condition cues & Posture and interaction details & Coat markings and individual morphology\\
\addlinespace
Pose and anatomy & Region alignment; joint locations alone may omit condition cues & Body configuration and its changes & Possible alignment aid; limited identity texture\\
\addlinespace
Viewpoint & Visibility of condition-related regions & Interpretation of projected posture & Correspondence across views\\
\addlinespace
Temporal context & Repeated observations of condition & Motion, duration, and transitions & Appearance continuity across observations\\
\addlinespace
Scene context & Often secondary to body evidence & Potentially useful near feed or water resources & Possible background-dependent matching\\
\bottomrule
\end{tabularx}
\end{table}

\subsection{Temporal Representations for Activity Recognition}

Temporal modeling adds information about change over time. The exact information depends on the representation. A sequence of RGB features keeps appearance and posture. A sequence of keypoints gives a smaller representation of body configuration. A video model can learn spatial and temporal patterns together. These choices keep different information and require different amounts of computation and annotation.

SlowFast networks make the difference between appearance and motion explicit by using two paths with different frame rates~\cite{feichtenhofer2019slowfast}. The slow path focuses more on spatial information, while the fast path captures quicker changes. SlowFast is also used as a benchmark in CVB~\cite{zia2023cvb}, connecting a general video-recognition method to cattle behavior. This gives a stronger reason for using temporal models than simply saying that more frames must be better.

Temporal convolution is another option. Bai et al.~\cite{bai2018tcn} compare convolutional and recurrent sequence models and show that temporal convolutions can work well across several sequence tasks. These models use temporal structure without requiring a recurrent hidden state at every step. For cattle behavior, this supports using a lightweight temporal module over image features when several visual representations need to be compared under limited compute.

A more structural option is to model body landmarks across time. ST-GCN, introduced by Yan et al.~\cite{yan2018stgcn}, represents a moving human skeleton as a graph with spatial and temporal links. The main idea is relevant to animal behavior because it models relationships between body parts instead of treating every coordinate separately. However, using it for cattle still depends on having suitable landmark definitions and reliable pose estimates. Pose sequences should therefore be treated as a separate representation type, not as an automatic replacement for RGB video.

These methods show that different temporal inputs answer different questions. Averaging several frame features gives the model more observed appearance, while a temporal model can also use the order of those frames. Pose-based features focus on body configuration, while dense video keeps richer motion and scene information. The best choice depends on the activity and the sampling interval. Therefore, this thesis uses temporal modeling as a focused representation question instead of simply increasing model size.

\subsection{Shortcut Learning, Domain Shift, and Evaluation}

Good performance in one familiar setting does not guarantee that the learned representation will work somewhere else. Geirhos et al.~\cite{geirhos2020shortcut} show that deep models can use easy correlations that fail when the environment changes. This is especially important when the object and its background appear together many times. In cattle images, the same pen, lighting, and camera position may appear repeatedly, making it easy for a model to learn scene-specific shortcuts.

Background dependence can be tested directly. Xiao et al.~\cite{xiao2021background} change foreground and background information to study their effect on image recognition. Their work shows why evaluation should go beyond testing on an unchanged dataset. For cattle-centered learning, the question is whether the model still depends on the cow when irrelevant parts of the scene change. Foreground processing is one way to test this, while changes in camera, setting, or dataset provide other robustness tests.

Animal-recognition research shows that this problem also appears in real-world settings. Beery et al.~\cite{beery2018terra} report a large performance gap between familiar and unfamiliar camera-trap locations. Wildlife camera traps are different from livestock cameras, but both can contain repeated fixed backgrounds. WILDS extends this idea with benchmarks that include natural distribution shifts across wildlife and other applications~\cite{koh2021wilds}. These studies support evaluating the conditions under which a model should generalize instead of treating every held-out image as equally difficult.

Domain shift can involve more than the background. It may include new cows, different breeds, changes in body appearance, new camera heights, or even different label definitions. Re-ID studies test some of these changes through multi-camera and longitudinal settings~\cite{yu2025multicam,zhang2026beca}. BCS studies may use rear, side, or depth views, while behavior datasets may use different activity definitions and environments. External evaluation therefore needs to consider both the image conditions and whether the labels mean the same thing.

Evaluation samples also need to be properly separated. Kapoor and Narayanan~\cite{kapoor2023leakage} show that data leakage can make machine-learning results look better than they really are. In video data, two different image files may still come from the same continuous event. A random image split can therefore place nearly identical situations on both sides of the train--test boundary. Grouping by the correct unit is safer. The correct unit depends on the claim: unseen-cow evaluation needs reliable cow IDs, while recording-based grouping protects against a different type of overlap.

This means that evaluation design is part of the representation question. A mask may appear useful under one split because the model is solving an easier identity- or location-related problem. Using several datasets may add visual variety while still leaving one class tied to one source. Training, comparison, and evaluation therefore need to match the claim being made. This gives a stronger basis for asking whether cattle-centered information improves transfer rather than only changing performance inside one familiar dataset.

\subsection{Multi-Task Learning: Selective Sharing and Negative Transfer}

MTL provides a way to combine the three task-specific models into one system. The simplest design shares an encoder and uses a separate head for each task. This can save computation and encourage common features, but it also forces several objectives to update the same representation. Crawshaw's review~\cite{crawshaw2020mtl} shows that architecture, task relationships, and optimization all affect MTL performance. Therefore, sharing needs a stronger reason than simply saying that all three inputs contain cattle.

Some MTL methods focus on balancing the task updates. Kendall et al.~\cite{kendall2018multi} use task uncertainty to weight losses with different properties. GradNorm changes task weights using gradient information and relative learning progress~\cite{chen2018gradnorm}. Both methods try to stop one task from dominating joint training because of loss scale or learning speed. They control how strongly tasks update the network, but they do not decide which visual features should be shared.

PCGrad addresses a different part of the problem by changing gradients when task updates conflict~\cite{yu2020pcgrad}. This makes gradient interaction part of the optimization method. Changing the size of an update is different from changing a conflicting direction, so PCGrad is not the same as loss weighting. However, both methods still work inside a chosen architecture. They help study optimization as one possible cause of negative transfer, but they do not prove that optimization is the only cause.

Other methods change how features are shared. Cross-stitch networks learn how to combine activations from task-specific networks~\cite{misra2016crossstitch}. The Multi-Task Attention Network by Liu et al.~\cite{liu2019mtan} uses task-specific attention to select information from a shared feature pool. These methods avoid forcing every task to use exactly the same feature path. This is useful when tasks overlap but still need different information, as with cattle morphology, activity, and identity.

Task grouping can also be tested instead of decided in advance. Standley et al.~\cite{standley2020tasks} study which tasks should be trained together under limited resources. Their work shows that task relationships cannot be decided only from task names or descriptions. Two tasks may share early visual features but need separate processing later, or one grouping may help one task while hurting another. Strong single-task references are therefore necessary before judging a shared model.

Animal-focused work shows both the value and the limits of joint learning. SAAM-VetNet combines disease detection and severity grading using an attention-based MTL framework~\cite{attri2025saamvetnet}. CATR-Net combines cattle-face segmentation and recognition~\cite{feng2025catr}. These studies show that MTL and perception-guided recognition are already useful in animal vision. However, their task relationships are different from those in this thesis. Segmentation and face recognition use the same visible target, while disease type and severity may use similar local signs. BCS, behavior, and Re-ID need more different kinds of information.

For example, Re-ID needs stable differences between cows. Behavior recognition must still respond when the same cow changes activity. BCS needs morphology but should not replace condition information with cow-specific appearance. A shared encoder can support all three only if it keeps these differences. This is why the thesis compares hard sharing with task-conditioned or partly private processing.

Using separate datasets for the tasks creates another issue. Joint training sees differences in camera view, image composition, and annotation style at the same time as it sees different task objectives. Some apparent task differences may therefore actually come from the datasets. A careful comparison must keep each task protocol fixed and examine every output separately. This makes it possible to test whether selective sharing helps the integrated model rather than simply giving it more parameters.

\begin{table}[tbp]
\centering
\caption[Approaches to managing multi-task interaction]{Complementary approaches to managing multi-task interaction. The methods address different aspects of joint learning and therefore support different interpretations of an improvement.}
\label{tab:ch2_sharing_methods}
\small
\begin{tabularx}{\textwidth}{@{}>{\raggedright\arraybackslash}p{3.0cm}YY@{}}
\toprule
\textbf{Approach} & \textbf{What changes} & \textbf{Question addressed}\\
\midrule
Uncertainty weighting~\cite{kendall2018multi} & Relative contributions of task losses & How should losses with different uncertainty be balanced?\\
\addlinespace
GradNorm~\cite{chen2018gradnorm} & Adaptive task weights using gradients & How should task learning rates be balanced?\\
\addlinespace
PCGrad~\cite{yu2020pcgrad} & Interaction of conflicting updates & Can harmful gradient interference be reduced?\\
\addlinespace
Cross-stitch sharing~\cite{misra2016crossstitch} & Exchange between task-specific features & Which feature combinations should be shared?\\
\addlinespace
MTAN~\cite{liu2019mtan} & Task-dependent selection of common features & Which shared information is useful to each output?\\
\bottomrule
\end{tabularx}
\end{table}

\section{Summary of Key Findings}

The literature shows that the research problem is more specific than simply combining three cattle classifiers. BCS, behavior, and Re-ID need different visual information. Anatomy-based scoring, temporal behavior models, and metric-learning Re-ID each solve part of the problem, but they focus on different properties of the cow. A unified framework therefore needs to keep these differences in mind.

\subsection{A Common Animal, but Different Information Requirements}

BCS research mainly focuses on anatomical coverage, body geometry, and the ordered BCS target. Behavior research adds movement, duration, and interaction context. Re-ID research focuses on keeping individual appearance consistent across different observations. These goals can support each other, but they can also conflict. Removing individual texture may help BCS avoid identity shortcuts, while Re-ID needs that texture. Scene context may help recognize feeding, while the same context can create background shortcuts for identity.

Previous research provides many ways to localize cattle, describe their body structure, and combine perception with recognition. The research gap is understanding how these choices affect BCS, behavior, and Re-ID differently, and how task-appropriate cattle-centered representations compare with generic RGB features when the tasks are brought into one shared system.

\subsection{Representation Quality Must Be Connected to Robust Evaluation}

Perception models can provide masks, landmarks, and viewpoint information, but producing an output is not enough. The representation must still keep the information needed by the task when relevant conditions change. For BCS, this includes visible anatomy and acquisition conditions. For behavior, it includes posture, motion, and scene changes. For Re-ID, it includes different views of the same cow across place or time.

This creates a second research gap: representation changes need to be tested under evaluation conditions that match the expected type of change. Foreground-focused models should be judged together with cross-setting and class-sensitive results, not only with one pooled score. Temporal models should also be evaluated in relation to how the observations are sampled over time. The thesis is therefore organized around the connection between visual information, task needs, and generalization conditions.

\subsection{Selective Sharing as the Basis of the Unified Framework}

MTL research separates two questions: how task updates are balanced, and what information is shared between tasks. The preliminary investigations in this thesis make this distinction important for cattle monitoring. Instead of assuming that heavy sharing is best, the study first uses single-task references to understand the representation needs and then compares different ways of combining the tasks.

The proposed research direction is therefore a task-conditioned cattle-centered framework. It aims to keep useful common cattle features while allowing morphology, activity, and identity to receive different emphasis. The main comparison is against generic RGB references and hard sharing, with attention to task-level performance and negative transfer. The goal is to understand which features can be shared, which need task-specific processing, and how those choices affect the unified system.

This literature review provides the basis for the remaining chapters. Chapter 3 describes the requirements and practical constraints. Chapter 4 explains the data protocols, representations, and model comparisons. Chapter 5 uses the measured results to return to the research questions.
